\documentclass{article}
\usepackage{amsmath}

\usepackage[utf8]{inputenc}     % pdfLaTeX
\usepackage[T1]{fontenc}
\usepackage[magyar]{babel}

\title{A Magyar nép}
\author{Lengyel Zoltán Péter, Rovenszky Robin}
\date{Legutóbb frissítve: 2026. Március 20.}

\newtheorem{definition}{Definíció}

\begin{document}

\maketitle
\tableofcontents
\newpage

\section{A Magyar nép őstörténete}
\begin{definition}[Őstörténet]
    Egy nép származása, eredettörténete.
\end{definition}

\begin{definition}[Nép]
    Közös kultúra, közös hagyomány, közös nyelv! Közös történelem, közös múlt, közös azonosságtudat!
\end{definition}
Az őstörténet forrásai.
\begin{itemize}
    \item Hagyományok, szokások
    \item Tárgyi emlékek; írásos emlékek; történetek.
    
    \begin{itemize}
        \item Rege
        \item Monda
        \item Geszta; pl. Geszta Hungarorum: A Magyarok őstörténete
    \end{itemize}
\end{itemize}
Anonimusz.
\begin{itemize}
    \item Nevét nem tudjuk
    \item Valószínűleg III vagy IV Béla írnoka volt 
\end{itemize}
Nomád életmód.
\begin{itemize}
    \item Halászó, vadászó életmód
    \item Ehhez kapcsolódó finnugor szavak
\end{itemize}
Magyar alapszókincs összetétele.
\begin{itemize}
    \item Szláv
    \item Germán
    \item Román
    \item Török
    \item Görög
\end{itemize}

\end{document}